% §8 Concluding remarks --- paper-only, written 2026-09-15 for module B's scope (blueprint
% chapters 8, 9, 10, 11 and 14; ADR-0006); rewritten 2026-09-18 at the author's review, on the
% hub's writing standard (writing-style-and-ai-use.md: claim first, one main clause and at most
% one subordinate, lists where the content is a list, no closers). On the line paper's model:
% the summary, the scope, one methodological remark on what the formalization changed in the
% printed proofs (each item is a blueprint CHANGED marker in chapters 8-11), and the open
% questions, stated as open and not as promised (hub RELEASES.md rule 1). Module C's subjects
% are named as the next module's without a statement number. No blueprint node is shared here.

\section{Concluding remarks}\label{sec:conclusions}

The line paper \sscite{fagerstrom2026spatial} characterized the spatial scale spaces on the
line under the hemigroup axioms: their kernels are the symmetric self-decomposable laws, each given by a
Gaussian coefficient and a nonincreasing displacement profile. This paper described that
class, and its results are of five kinds.

\emph{The shape of the class.} The admissible exponents form a cone. Its extreme rays are the
Gaussian ray and the exponents $\Cin(\tau\,\cdot)$ of the unit-step profiles, and every member
is a unique superposition of them (Proposition~\ref{prop:choquet-cone}). The kernel of a
member with no Gaussian part is bounded at the origin exactly when the value of its profile
there exceeds $1$. The
$\Cin$ rays sit at that threshold, and so does the \Matern{} member at $\gamma = \tfrac12$
(Corollaries~\ref{cor:origin-boundedness} and~\ref{cor:origin-smooth}).

\emph{The relation to the causal theory.} Bochner's subordination maps the causal theory of
\sscite{fagerstrom2026hemigroup} into the spatial one, family to family and corner to corner
(Lemma~\ref{lem:bridge-exponents}, Proposition~\ref{prop:bridge-families}). The map is
injective (Proposition~\ref{prop:bridge-strictness}(1)), and on the Thorin subclasses it is a
bijection (Proposition~\ref{prop:thorin-subclass}). It is not onto. It misses every member
with no Gaussian part whose profile is a lattice step, the $\Cin$ rays among them
(Proposition~\ref{prop:bridge-strictness}). Its
image is the intersection of a chain of classes indexed by dimension, the exponents whose
radial extension is self-decomposable in dimension $d$, and the first two inclusions of the
chain are strict (Proposition~\ref{prop:dimension-filtration}, proved in prose on cited facts
about laws in dimension $d$, and Lemma~\ref{lem:filtration-strictness}, not machine-checked).

\emph{The named members.} The \Matern{} family is the single-atom member of the Thorin
subclass with no Gaussian part. Its members of integer index are exactly the families
realized at every ladder of scales by a cascade of identical first-order forward--backward
sections, and every member has every moment finite and exponential tails
(Theorem~\ref{thm:matern},
Corollary~\ref{cor:matern-realization}). The Student-t family is the image of the causal Bessel
corner, with polynomial tails of a prescribed order (Proposition~\ref{prop:student-t}). The
symmetric stable family is the semigroup corner, and the heavy tails of its non-Gaussian
members come from the semigroup requirement (Proposition~\ref{prop:stable-family}).

\emph{The evolution.} Every member satisfies an evolution equation in scale whose generator is
a Laplacian plus a superposition of symmetric second differences
(Proposition~\ref{prop:scale-evolution}). The generators of the named members are a
differential operator, a fractional one, a bounded integral operator, mixtures of such, and
for the unit-step members a single second difference
(Proposition~\ref{prop:corner-generators}).

\emph{The algorithm.} The cascade is the algorithm. Over any finite ladder of scales the
cascade is exact at every knot, for a signal on the line, the sampling of the signal being a
separate approximation. At integer index the \Matern{} stages are forward--backward
first-order recursive sections. The members of the Thorin subclass with finitely many ranges
and integer weights have such stages, and every member of the subclass is a limit of members
with finitely many ranges. The restriction to integer weights is a real one: the limits of
the families with such stages are the members whose Thorin measure consists of atoms of even
integer mass, finitely many in every bounded set, possibly with a Gaussian part. Every such
limit has an exponential moment, so the stable and the Student-t members are not among them
(Proposition~\ref{prop:rational-closure}). The numerical run of \S\ref{sec:numerical-run} finds them at a fixed distance from the Poisson
scale space.

\emph{Scope.} The paper delivers the geometry of the class, the bridge, the corners, the
generator form and the implementation, on the line and in the $L^1$ signal model. It makes no
empirical claim about natural signals. Locality of the scale evolution and of the scale space,
the non-creation and non-enhancement of local extrema, and the scale-space jet are the subject
of the next module. Nothing here depends on them. Dimension $d$, and families that are
covariant only under a discrete group of dilations, remain outside, as in the line paper.

\emph{Open questions.} Five are left.
\begin{enumerate}
\item How large the complement of the slice is. The lattice-step members with no Gaussian part
are outside it, and so
is every member with a compactly supported catalogue. No description of the complement is
given.
\item Which admissible families have rational stages. The families realized by first-order
sections are among them, and so are a family with a complex pole pair and a family with real
poles and zeros only (Remark~\ref{rem:rational-converse}), so the Thorin subclass does not
contain them all. Proposition~\ref{prop:rational-criterion} says which rational functions are
kernels from the origin. Which are stages between two positive scales, and how the others are
realized with positivity at every stage, is not attempted.
\item The formal status of three of the clauses that \S\ref{sec:machine-checked} lists as not
formalized. The closed form of the \Matern{} kernel, its tail, and the moment count of the
Student-t family are cited and not machine-checked
(Propositions~\ref{prop:matern-density} and~\ref{prop:student-t}(4)). A machine-checked closed
form would let the boundedness threshold be read from the kernel as well as from the profile.
\item The boundedness threshold in dimension $d$. On the line it is the value $1$ of the
profile at the origin, and the condition $\gamma > \tfrac12$ of the \Matern{} member is its
instance. Whether the threshold in $\RR^d$ is the dimension is a question for the theory in
dimension $d$, which is not begun here. If it is, an admissible aperture in the plane may be
unbounded at its own centre.
\item The hemigroup weakened once more. The dilatively stable processes of
\sscite{kern2015dilatively} and the semi-selfsimilar and semistable hemigroups of
\sscite{beckerkern2004random} describe what remains of covariance when the scaling group is
weakened or discretized. How much of the present geometry survives there is open.
\end{enumerate}
